% Part: applied-modal-logics
% Chapter: epistemic-logic
% Section: properties-accessibility

\documentclass[../../../include/open-logic-section]{subfiles}

\begin{document}

\olfileid{aml}{el}{acc}

\olsection{可达关系与认知原则}

基于我们在正规模态逻辑中对框架对应性的已有了解，我们可能会想看看在认知解释下特征公式是什么样子。
我们之前已经提到，认知逻辑通常是在 S5 中解释的。
所以，让我们看看各种认知原则是如何得到表达的，
并考虑它们如何对应于各种框架条件。

回顾正规模态逻辑可知，不同的模态公式刻画了可达关系的不同性质。
下表列出了其中几个对应于特定认知原则的性质。

\begin{table}[t]
    \begin{tabular}{| p{.48\textwidth} || p{.48\textwidth} |}
      \hline
      {\emph{若 $R$ 是……}} & {\emph{则……在~$\mModel{M}$ 中为真：}} \\
      \hline \hline

      & \hspace{2em} $\Knows (p \lif q) \lif (\Knows p \lif \Knows q)$
      \hfill \newline{(封闭性)} \\
      \hline
      \emph{自反的}：$\forall w Rww$
      & $\Knows p \lif p$ \hfill \newline{(真实性)} \\
      \hline
      \emph{传递的}：\newline
      $\forall u \forall v \forall w ((Ruv \land Rvw) \lif Ruw)$ &
      $\Knows p \lif \Knows \Knows p$ \hfill
           \newline{(正内省)} \\
      \hline
      \emph{Euclid的}：\newline
      $\forall w \forall u \forall v ((Rwu \land Rwv) \lif Ruv)$ &
      $\lnot \Knows p \lif \Knows \neg \Knows p$ \hfill
        \newline{(负内省)} \\
      \hline
    \end{tabular}
    \caption{四个认知原则。}
    \ollabel{tab:four}
  \end{table}

对应 $T$ 公理的“真实性”原则，通常被视为这些原则中最无争议的一个，因为它表达的主张是：如果一个公式被知道，那么它必定为真。
而“封闭性”以及“正内省”和“负内省”原则则争议大得多。

“封闭性”对应 $K$ 公理，它代表了主体的知识在蕴涵下封闭这一观点。
在某些情况下，这似乎是合理的。
例如，我可能知道：如果我在维多利亚市，那我就在温哥华岛上。
在不考虑离奇的怀疑论情景的前提下，我确实知道自己在维多利亚市，而这似乎也意味着我知道自己在温哥华岛上。
所以在这个案例中，我的知识的逻辑封闭性显得相对直观。
另一方面，我们并不总会推演自己知识的所有推论，因此在其他情况下，这可能导致不那么直观的结果。

“正内省”有时被称为 KK 原则，常被表述为：如果我知道某事，那么我知道我知道它。
它是 $4$ 公理的认知对应物。
相应地，“负内省”的表述是：如果我\emph{不知道}某事，那么我知道我不知道它，这是 $5$ 公理的对应物。
这两者似乎都存在相对常见的反例，比如这样一些情况：我实际上知道某件事，却不确信自己是否知道它。

\end{document}